\section*{Terminal-state sufficiency for rank-allocation sequences}
\label{sec:terminal-state-sufficiency}

\paragraph{Scope.}
This is a structural proposition about the evaluation model used by the
adaptive tensor-network campaign, not a statement about projected multilinear
trees in general. It is elementary. It is recorded because it is load-bearing:
it is what licenses replacing a search over action sequences by an enumeration
over terminal rank vectors, and the whole $m=8$ terminal-oracle calibration
rests on it.

\paragraph{Setting.}
Fix a finite ordered rooted tree $T$ with internal vertices carrying
multilinear laws $\mu_v$, a finite leaf batch $Z$, and a family of orthogonal
projectors $\{P_v^{(\cdot)}\}$ in which $P_v^{(\rho)}$ denotes projection onto
the leading $\rho$ directions of a basis fixed in advance. Write
$\mathbf r \in \mathbb{Z}_{\ge0}^{V}$ for a rank assignment and
\[
  \Phi(Z,\mathbf r) \;:=\; \bigl\|\,F_r(Z) - \widetilde R_r(Z,\mathbf r)\,\bigr\|
\]
for the root discrepancy, where $F$ is the ambient evaluation and
$\widetilde R$ the recursively projected one.

\paragraph{Hypotheses.}
\begin{enumerate}
\item[(H1)] \emph{Fixed projector bases.} Each $P_v^{(\rho)}$ is determined by a
basis fitted once, from the ambient pass, independently of $\mathbf r$.
\item[(H2)] \emph{Purely rank-indexed evaluation.} $\widetilde R$ is computed by
a recursion over $T$ that reads only $(Z, \mathbf r)$ and the laws, and retains
no state between invocations.
\end{enumerate}

\paragraph{Proposition (terminal-state sufficiency).}
Under (H1)--(H2), for any two admissible action sequences
$a_{0:T-1}$ and $a'_{0:T-1}$ with
\[
  \mathbf r_0 + \sum_{t=0}^{T-1} a_t \;=\; \mathbf r_0 + \sum_{t=0}^{T-1} a'_t,
\]
the terminal objectives coincide:
\[
  \boxed{\;\Phi\bigl(Z,\mathbf r_T\bigr) = \Phi\bigl(Z,\mathbf r_T'\bigr).\;}
\]

\emph{Proof.} By (H1)--(H2) the map $\mathbf r \mapsto \Phi(Z,\mathbf r)$ is a
function --- the evaluation depends on the sequence only through the rank
vector it produces. Equal terminal rank vectors therefore give equal
objectives. $\square$

\paragraph{Corollary (enumeration validity).}
Let each round add one unit to exactly $b$ distinct coordinates, over $T$
rounds. Then the reachable terminal increments are exactly
\[
  \mathcal X_{T,b}
  \;=\;
  \Bigl\{\,x \in \{0,\dots,T\}^m \;:\; \textstyle\sum_i x_i = Tb \,\Bigr\},
\]
the constraint $x_i \le T$ and the total being precisely the Gale--Ryser
condition for decomposing $x$ into $T$ binary vectors of weight $b$. By the
proposition, minimizing the terminal objective over policies is equivalent to
\[
  \boxed{\;\min_{x \in \mathcal X_{T,b}} \Phi\bigl(Z, \mathbf r_0 + x\bigr).\;}
\]
For $m=8$, $b=3$, $T=6$ this replaces $\binom{8}{3}^6 = 3.08\times10^{10}$
action sequences by $|\mathcal X_{T,b}| = 235{,}348$ terminal profiles.

\paragraph{Boundary.}
The proposition fails if the projector bases are refitted from the reduced
pipeline at each round, since then $P_v$ would depend on the history through
the values it was fitted to; (H1) excludes exactly that. It also says nothing
about the difficulty of the optimization: the sequential problem remains
combinatorial because the per-round budget constrains which terminal profiles
remain reachable from a given partial state. What it removes is hysteresis, not
planning.

\paragraph{Implementation control.}
\texttt{run\_terminal\_oracle\_m8.py} tests the proposition numerically on 20
random terminal profiles per instance, each realized by 4 distinct
decompositions, and observes a maximum spread of $0.000\mathrm{e}{+}00$. Under
the proposition this is a test of the implementation, not evidence for the
statement.
